
% Section 5 rewrite for Paper 2
% Assumes Proposition 4.2 / full spanning statement is in place

\section{Minimality of Degree $4$ and the First Recovery}\label{sec:minimality}

The aim of this section is to show that, once the basis is enlarged from
\[
V_d^{(6)}=\Span_{\mathbb Q}\{n^kM_a(n):0\le k\le d,\ 1\le a\le 6\}
\]
to
\[
V_d^{(7)}=\Span_{\mathbb Q}\{n^kM_a(n):0\le k\le d,\ 1\le a\le 7\},
\]
the first new prime-vanishing direction involving the obstruction sector occurs at degree $4$. The argument splits naturally into two parts:

\begin{enumerate}
    \item a structural description of the obstruction space $O_d$, and
    \item exact rational verification of the minimal degree at which a nontrivial null vector appears outside the previously generated span.
\end{enumerate}

The structural description is provided by Proposition~\ref{prop:Od-full-structure}, which shows that
\[
O_d=\Span_{\mathbb Q}\{[\tau],[p\tau],\dots,[p^{d+1}\tau]\},
\qquad
\dim O_d=d+2.
\]

Thus the enlargement by $M_7$ does not remove the cusp obstruction; it enlarges the space in which that obstruction may be cancelled.

\subsection{The cusp cancellation constraint}

Let $R(n)\in V_d^{(7)}$ be a prime-vanishing expression. Passing to prime restrictions and then modulo $E_d$, only the cusp sector survives. By Proposition~\ref{prop:Od-full-structure}, the obstruction sector is spanned by
\[
[\tau],[p\tau],\dots,[p^{d+1}\tau].
\]

Write the $M_6$- and $M_7$-parts of $R$ as
\[
Q_6(n)M_6(n)+Q_7(n)M_7(n),
\]
where $Q_6,Q_7\in\mathbb Q[n]$ have degree at most $d$. Using the prime-level decompositions from Section~3,
\[
M_6(p)=f_6(p)+c_\tau\tau(p),
\]
\[
M_7(p)=f_7(p)+\alpha\tau(p)+\beta p\tau(p),
\qquad \beta\neq 0,
\]
and discarding the Eisenstein terms modulo $E_d$, the cusp contribution of $R(p)$ becomes
\[
\bigl(c_\tau Q_6(p)+\alpha Q_7(p)+\beta p\,Q_7(p)\bigr)\tau(p).
\]
Therefore a necessary and sufficient condition for the cusp obstruction to vanish modulo $E_d$ is
\begin{equation}\label{eq:cusp-constraint}
c_\tau Q_6(p)+(\alpha+\beta p)Q_7(p)=0
\end{equation}
as a polynomial identity in $p$.

Equation~\eqref{eq:cusp-constraint} is the structural cancellation constraint governing the recovery problem.

\subsection{Degrees $0$ and $1$}

At degree $d=0$, the obstruction space $O_0$ is two-dimensional by Proposition~\ref{prop:Od-full-structure}, with basis $[\tau]$ and $[p\tau]$. The available freedom in the $M_6$- and $M_7$-sectors is only two scalar coefficients, and Equation~\eqref{eq:cusp-constraint} forces these coefficients to vanish unless the entire cusp contribution is trivial. Thus no new recovered prime-vanishing direction appears at degree $0$.

At degree $d=1$, the same structural obstruction persists. The classes $[\tau],[p\tau],[p^2\tau]$ are linearly independent in $O_1$, while the cancellation constraint \eqref{eq:cusp-constraint} still forces any potential prime-vanishing combination to annihilate a three-dimensional obstruction sector using insufficient freedom. Hence no new recovered direction appears at degree $1$.

The absence of recovery at degrees $0$ and $1$ is therefore structural rather than computational.

\subsection{Degrees $2$ and $3$: exact verification}

For degrees $d=2$ and $d=3$, the structural cusp constraint remains necessary but no longer by itself rules out cancellation. In these cases we use exact rational computation.

\begin{proposition}\label{prop:no-recovery-d23}
For $d=2$ and $d=3$, the exact prime-evaluation matrices for the enlarged basis $V_d^{(7)}$ contain no new prime-vanishing direction beyond the previously generated span.
\end{proposition}

\begin{proof}
This is verified by exact rational null-space computation on the prime-evaluation matrices using the first $95$ primes in $[2,500]$. In degree $2$, the nullity gain after adjoining $M_7$ is $0$. In degree $3$, the nullity gain is also $0$. Hence no new recovered direction appears in either degree.
\end{proof}

\subsection{Degree $4$: first recovery}

At degree $d=4$, the same exact rational computation exhibits the first nullity gain.

\begin{proposition}\label{prop:first-recovery-d4}
In degree $4$, the enlarged basis $V_4^{(7)}$ admits a new prime-vanishing direction not contained in the previously generated span. This first recovered direction is denoted $E_6$.
\end{proposition}

\begin{proof}
The exact prime-evaluation matrix for $V_4^{(7)}$ has nullity gain $1$ relative to the restricted basis, and composite-value span testing shows that the resulting null vector is independent of the span generated by the previously known expressions. This yields a new recovered direction, denoted $E_6$.
\end{proof}

Combining Propositions~\ref{prop:no-recovery-d23} and \ref{prop:first-recovery-d4}, we obtain the main minimality statement.

\begin{theorem}\label{thm:minimal-degree-4}
The first new recovered prime-vanishing direction in the enlarged basis $\{M_1,\dots,M_7\}$ appears at degree $4$. No such direction appears at degrees $0$, $1$, $2$, or $3$.
\end{theorem}

\begin{proof}
The cases $d=0,1$ are ruled out structurally by the obstruction-space analysis above. The cases $d=2,3$ are ruled out by Proposition~\ref{prop:no-recovery-d23}. Degree $4$ is the first degree at which Proposition~\ref{prop:first-recovery-d4} yields a new recovered direction.
\end{proof}

\subsection{Interpretation}

The point of Theorem~\ref{thm:minimal-degree-4} is that the appearance of $E_6$ is not accidental. The obstruction sector is already present at degree $0$, but the enlarged basis does not admit its first successful cancellation until degree $4$. In this sense, $E_6$ is the first genuine recovery after the weight-$12$ barrier.
